\section{Edge collisions for alphabets of size at least three}
\label{sec:qary-edge-poisson}

We now establish the Poisson limit for colliding pairs of edges when the alphabet has size at least three. Fix \(q\ge3\). As in Section~\ref{sec:classical-poisson}, set \(s=n^{1/3}\) and \(p_0=1-1/q\), and retain the constants \(A,v,R,t_*,D,\kappa_1\) defined there. Define
\[
 \Delta=\frac{R^3(3p_0-2)}{6vp_0}=\frac{q-3}{q-1}\log q,\qquad
 K=\kappa_1\sqrt n\exp\left(As+\frac{t_*}{s}
              +\frac{(9/(4R))\log n+y}{s^2}\right).
\]
For any fixed \(y\), the number of colliding unordered pairs of distinct edges converges in distribution to a Poisson random variable with mean
\[
 \frac{(q-1)^2}{8}(2\pi v)^{1/4}\exp(-D-\Delta-Ry).
\]
Consequently, the resolution probability converges to the exponential of minus this mean. Both landmark models yield the same limit: independent selection of each vertex with probability \(K/q^n\), and uniform selection of a subset of size equal to the nearest integer to \(K\).

To prove this limit, we adapt the vertex argument of Section~\ref{sec:classical-poisson} to edges. We first determine how edge distances alter the shell intensities and the scalar expansion, next verify the Fourier and determinant estimates on typical configurations, and finally bound the atypical configurations and pass to factorial moments.

\step{Shell intensities and the scalar expansion.}
We begin by comparing the distance distribution from a uniformly random landmark to a fixed edge with the corresponding distribution for a vertex, and computing the resulting shift in the scalar expansion. The distance from a random landmark to an edge (the edge radius) is the sum of a \(\mathrm{Bin}(n-1,p_0)\) variable, arising from the coordinates where the two endpoints of the edge agree, and an independent \(\mathrm{Bern}((q-2)/q)\) variable, arising from the single active coordinate where the landmark mismatches both endpoint letters. Dividing this convolution by the vertex shell mass yields the exact ratio \(2-r/(np_0)\).

In the rescaled radial coordinate \(x=(r-np_0)/s^2\), the resulting change to the log-intensity profile \(h_r\) is \(-x/(p_0s)-x^2/(2p_0^2s^2)+O(s^{-3})\) on the moderate window. In the clipped half-sum, the odd linear term integrates to zero over the interior, while the quadratic term contributes \(-R^3/(6p_0^2)\). At the boundary, the shift in \(j\) contributes a correction of \(c_3vR^3/p_0\), as given by the two-branch formula in Section~\ref{sec:classical-poisson}. Adding the interior and boundary contributions yields \(\Delta\).

The remaining parts of the scalar expansion are unchanged to leading order: the boundary density for the \(G\) sum through constant order still depends only on \(F_0,g\), which are unaffected, and the central deletion terms change by \(o(1)\). Replacing the vertex coefficient \(u_*=1/(4R)\) by \(9/(4R)\) in the scalar expansion therefore yields the single-pair collision atom
\[
 (1+o(1))\,q^{-2n}n^{-2}(2\pi v)^{1/4}
   \exp(-D-\Delta-Ry),
\]
once the relative Fourier and determinant estimates below are established.

\step{Typical Fourier estimates.}
Next, we verify that the Fourier localization, dense-box approximation, and small-row splitting from the vertex analysis carry over to typical collections of edge pairs. The strategy is to freeze the landmark letters on the active coordinates of the edges, reducing all edge distances to ordinary vertex distances on a common core up to bounded offsets.

Fix a list of edge pairs, and freeze the landmark letters on the set \(I\) of their active axes. Conditioned on these frozen letters, every edge distance is an ordinary Hamming distance on a common core of dimension \(n-|I|\) plus a bounded frozen offset. Call the list typical when all common-core tuple frequencies of its listed endpoints lie within \(O(n^{11/12})\) of their uniform means.

To adapt the singleton construction of Section~\ref{sec:classical-birthday}, fix any frozen pattern on \(I\), subtract its offsets from the radius constraints, and subtract the corresponding offset differences from the equality constraints. Because the extra residuals introduced by the offsets are bounded, the same unimodular moves inside linear-size core cells still solve all constraints. Fixing the frozen pattern on \(I\) costs only a constant factor in probability. Omitting the anchor radius constraint and imposing equality in every block likewise yields the polynomial-mass neutral word class, obtained by rounding the untilted allocation on the core and correcting the bounded residuals. The two-tilt bounds also hold on each frozen fiber, because bounded shifts leave their error orders unchanged. Thus singleton localization, dense-box approximation, and small-row splitting all apply with the edge shell intensities. These singleton and neutral classes also supply the hypotheses required for the relative fixed-size transfer of Section~\ref{sec:fixed-cardinality}, including the step that excludes repeated draws on full words.

\step{Determinant estimates and factorization.}
We now compute the Gaussian determinant for a typical edge pair and show that the joint determinant for multiple typical pairs factors asymptotically into the product of the individual block determinants. Write \(\pi\) for the common edge distance marginal and \(w=(\sqrt{\pi_r})_r\) for the unit vector of square-rooted shell probabilities. For a frozen pattern \(z\), the shift map \(T_{e,z}f(t)=f(t+a_e(z))\) from the edge radial space to the core radial space has operator norm at most \(q^{|I|/2}\), because \(\pi\) is a mixture over the frozen patterns. For two typical distinct core centers, the centered core joint operator has trace norm \(o(1)\) by the radial Krawtchouk estimate of Section~\ref{sec:classical-poisson}; conjugating by the bounded maps \(T_{e,z}\) preserves this \(o(1)\) trace-norm bound.

We next control the conditional mean density \(g_{e,z}\) of the edge distance given \(z\), taken relative to \(\pi\). Because this density arises from conditioning on a frozen pattern on \(I\), it is bounded by \(q^{|I|}\). Moreover, bounded shifts of the core binomial distribution have total variation distance \(O(n^{-1/2})\) from the unshifted distribution by unimodality and telescoping. Combining this total variation bound with the uniform bound on \(g_{e,z}\) gives \(\|g_{e,z}-1\|_{L^2(\pi)}^2=O(n^{-1/2})\). Since the densities \(g_{e,z}\) average to one over \(z\), the conditional mean outer products average to the constant projection plus an operator of trace norm \(o(1)\).

In whitened coordinates, the full joint operator therefore takes the form \(ww^T+F\), where \(Fw=F^Tw=0\) and \(\|F\|_1=o(1)\). Here \(F\) need not be symmetric. For a discrepancy between two edges, dividing the covariance matrix by the marginal diagonal scale \(2\lambda_r\) gives \(I-ww^T-(F+F^T)/2\). Applying the isometric compression of Section~\ref{sec:classical-poisson} yields the determinant \((1+o(1))\pi_c\) after omitting one central row and restricting to dense rows. For distinct discrepancy blocks, the constant projections cancel in their cross-covariance, and the same isometries show that the cross-covariance has trace norm \(o(1)\) after whitening the individual covariance blocks. Hence the ratio of the joint determinant to the product of the individual determinants tends to one. The small-row factors likewise factor relatively, by the two-tilt and singleton estimates.

\step{Atypical configurations and counting.}
It remains to bound the atypical configurations and sum the typical contributions to obtain the factorial moments. We adapt the forest estimates of Section~\ref{sec:classical-poisson} to the common core by adding one rule for choosing the frozen landmark letters on the active coordinates.

For an upper-layer anchor, choose the frozen letters to mismatch each anchor factor; for a lower-layer anchor, choose them to match each anchor factor. Because every active factor of an edge is either a singleton or a two-letter set, both a matching choice and a mismatching choice exist for \(q\ge3\). Under this rule, every comparator's frozen offset is weakly inward. To obtain strictness against a specified close comparator, we consider two cases. If the comparator has a different core word from the anchor, the core rounding construction already gives strictness. If its core word equals the anchor's, the two edges are still distinct, so their frozen product sets are distinct subsets of equal cardinality. Thus some comparator factor contains a letter outside the corresponding anchor factor, and some anchor factor contains a letter outside the corresponding comparator factor; otherwise one product set would contain the other. Choosing the first letter for an upper-layer anchor and the second for a lower-layer anchor makes the frozen offset strictly inward. Core rounding, restoration in an agreement cell, and a sufficiently wide fixed central exclusion strip then give the same leading-layer independence and inverse-polynomial relative shell mass as for vertices.

We now apply these forest bounds to eliminate atypical pairs and tuples. For a single close pair, nominate its other edge as the close comparator; the resulting single complete family gives an atom bound of \(\exp[-2Ln+O(n^{2/3}\log n)]\). At a fixed, sufficiently small cluster threshold \(\delta\), the number of such close pairs is at most \(\exp[(L+b_q(\delta))n+O(\log n)]\), so their total contribution vanishes. For a well-separated atypical pair, the forest gives the same atom bound, while the nonuniform core type frequencies reduce the count by a fraction \(\exp[-\Omega(n^{5/6})]\). Together these bounds remove all atypical pairs, including pairs with identical cores and differing frozen factors.

For factorial tuples with \(b\) clusters in \(c\) connected components, the same forest produces \(b-c\) complete families. In the case \(b=2c\) where some component has two distinct links, the anchors of those links differ on at least one side, so the strict comparison rule produces an extra complete family. The unchanged entropy comparison eliminates both of these clustered cases, and any remaining nontypical tuples incur the same type-frequency penalty. Choosing the active axes and frozen endpoint letters costs only polynomial factors, and all of these negligibility bounds absorb the \(O(\sqrt K)\) penalty from conditioning on a fixed landmark count.

\step{Counting and factorial moments.}
There are \(n(q-1)q^n/2\) edges in total. Typical unordered pairs have asymptotically the full unordered pair count, and typical ordered factorial tuples have asymptotically its corresponding powers: for any choice of active axes, the free core endpoints have uniform independent coordinate tuples before atypical or repeated pairs are discarded. Parallel pairs (sharing the same active axis) account for an \(O(1/n)\) proportion of all pairs and do not require a separate typical atom. Combining the relative factorization with the single-pair atom therefore shows that every fixed factorial moment converges to the corresponding power of the displayed Poisson mean. Since atypical single pairs vanish in probability, this proves the asserted Poisson and resolution probability limits.
